% Part: lambda-calculus
% Chapter: lambda-definability
% Section: pairs

\documentclass[../../../include/open-logic-section]{subfiles}

\begin{document}

\olfileid[hi]{lam}{ldf}{pai}
\olsection{युग्म और पूर्ववर्ती}

\begin{defn}
$M$ और $N$ का युग्म (जिसे $\tuple{M,N}$ लिखते हैं) इस प्रकार परिभाषित है:
\[
\tuple{M,N} \ident \lambd[f][fMN].
\]
\end{defn}
  
सहज रूप से यह ऐसा फलन है जो किसी फलन को स्वीकार करके उसे युग्म के दोनों
अवयवों पर लागू करता है। इसी विचार से हमारे पास यह निर्माता है, जो दो पद
लेकर उन्हें समाहित करने वाला युग्म लौटाता है:
\[
  \fn{Pair} \ident \lambd[mn][\lambd[f][fmn]]
\]
कोई युग्म दिए होने पर हम उसके अवयव भी पुनः प्राप्त करना चाहते हैं। इसके
लिए दो अभिगम फलन चाहिए, जो युग्म को तर्क के रूप में स्वीकार करके उसका पहला
या दूसरा अवयव लौटाते हैं:
\begin{align*}
  \fn{Fst} & \ident \lambd[p][p(\lambd[mn][m])]\\
  \fn{Snd} & \ident \lambd[p][p(\lambd[mn][n])]
\end{align*}

\begin{prob}
  समझाएँ कि अभिगम फलन $\fn{Fst}$ और $\fn{Snd}$ क्यों काम करते हैं।
\end{prob}

अब युग्मों से हम पूर्ववर्ती फलन को !!{lambda define} कर सकते हैं:
\[
  \fn{Pred} \ident \lambd[n][\fn{Fst}(n (\lambd[p][\tuple{\fn{Snd}\, {p}, \fn{Succ}(\fn{Snd}\, {p})}]) \tuple{\num 0, \num 0})]
\]
स्मरण करें कि $\num n\, f x$ अपचयित होकर $f^{n}(x)$ बनता है; इस स्थिति
में $f$ ऐसा फलन है जो कोई युग्म $p$ स्वीकार करके ऐसा नया युग्म लौटाता है
जिसमें $p$ का दूसरा घटक और उस दूसरे घटक का उत्तरवर्ती होता है; $x$ युग्म
$\tuple{0,0}$ है। अतः $n=0$ के लिए परिणाम $\tuple{0,0}$ और अन्यथा
$\tuple{\num{n-1}, \num n}$ है। फिर $\fn{Pred}$ परिणाम का पहला घटक लौटाता
है।

व्यवकलन को $a$ पर $\fn{Pred}$ को $b$ बार लागू करके परिभाषित किया जा सकता है:
\[
\fn{Sub} \ident \lambd[ab][b \fn{Pred}\, a].
\]
    
\end{document}
